\section{Introduction}

The characteristic initial value problem is a natural formulation of the
Einstein equations when radiation, null focusing, and the formation of
trapped surfaces are central.  Instead of prescribing data on a spacelike
hypersurface, one prescribes compatible data on two intersecting null
hypersurfaces and evolves their future domain of dependence.  Local
well-posedness for this problem is known under suitable regularity
assumptions~\cite{rendall-90,luk-12,cabet-16,hilditch-20}, but its direct
numerical realization becomes substantially more difficult once spherical
symmetry is removed.  The use of null hypersurfaces as an evolution framework
goes back to the Bondi--Sachs description of gravitational radiation and the
early analysis of characteristic data~\cite{bondi-62,sachs-civp-62}; its
development as a numerical method is surveyed in~\cite{winicour-12}.

Characteristic numerical relativity has progressed from double-null vacuum
evolutions with two Killing symmetries~\cite{corkill-stewart-83} and
spherically symmetric scalar collapse in null coordinates
~\cite{garfinkle-95,hamade-stewart-96} to axisymmetric vacuum evolution on
outgoing null cones~\cite{gomez-94} and generic angular data for accurate
Cauchy--characteristic waveform extraction~\cite{moxon-23}.  Numerical
relativity has also long evolved nonspherical spacetimes in spacelike
foliations.  Fully three-dimensional Cauchy calculations have studied
scalar-field collapse without symmetry assumptions~\cite{deppe-19}, and
axisymmetric pseudospectral calculations have reached strongly aspherical
regimes in which the center of collapse bifurcates~\cite{marouda-24}.  The
first null-coordinate simulations of nonspherical regular data collapsing
to a black hole were obtained in twist-free axisymmetry for the massless
scalar field~\cite{gundlach-24}; that work describes the accessible data as
moderately nonspherical and identifies larger deviations from spherical
symmetry, as well as vacuum collapse, as important challenges.  Thus the gap
addressed here is not an absence of nonspherical numerical solutions in
general.  It is the relative scarcity of direct characteristic evolutions
with strongly nonspherical, freely prescribed initial geometry in a fully
angular double-null gauge.

Let $(u,v,\theta^A)$ be double-null coordinates.  We write the spacetime
metric as
\begin{equation}\label{eq:introduction-double-null-metric}
  \bg=-2\Omega^2(\dd u\otimes\dd v+\dd v\otimes\dd u)
  +g_{AB}(\dd\theta^A-b^A\dd u)\otimes
   (\dd\theta^B-b^B\dd u).
\end{equation}
The metric is therefore represented by the section metric $g_{AB}$, the
lapse $\Omega$, and the angular shift $b$.  This formulation gives
considerable freedom to choose the geometry on the two initial null
hypersurfaces, subject to the characteristic constraints and corner
compatibility.  In particular, the nonspherical part of the initial metric
variation can be prescribed directly rather than produced only indirectly
from a matter perturbation or a spacelike constraint solve.  The fourth
experiment gives a quantitatively strong example: on its outgoing initial
hypersurface, where $\Omega=1$, the prescribed data satisfy
\begin{equation}\label{eq:introduction-strong-asymmetry}
  \norm{\tf\Lie_{\partial_v}g(-1,0.005)}_{L^2(S)}=7.50960,
  \qquad
  \norm{\tf\Lie_{\partial_v}g(-1,0.005)}_{L^\infty(S)}=2.53468.
\end{equation}
The perturbation is consequently a substantial geometric departure from
spherical symmetry, rather than a nonspherical profile whose effect on the
initial metric is numerically negligible.

This paper develops a first-order Picard iteration for both the Einstein
vacuum equations
\begin{equation}
  \Rc(\bg)=0
\end{equation}
and the Einstein--massless-scalar system
\begin{equation}
  \Rc(\bg)=\dd\phi\otimes\dd\phi,
  \qquad \Boxg\phi=0,
\end{equation}
in the gauge \eqref{eq:introduction-double-null-metric}.  The iteration is
organized by the geometry of the characteristic problem.  Data that are
free, data fixed by the null constraints, and quantities transported in the
two characteristic directions play different roles in each sweep.  The
resulting construction is therefore not obtained by treating all coordinate
components as an undifferentiated system of equations: its update order
follows the causal and geometric dependency structure of the double-null
equations.

Two difficulties are particularly important.  The first is to obtain an
executable form of the equations.  Analytic arguments can often group
lower-order terms schematically or pass between conventionally equivalent
forms, whereas numerical evolution requires every displayed sign and
numerical factor to be fixed consistently.  We address this by checking the
formulas against identities derived from the four-metric, by testing exact
solutions, and by evaluating curvature residuals independently of the
quantities used to construct a Picard sweep.  The second difficulty is the
iteration itself.  Its design requires identifying a closed hierarchy that
respects the characteristic constraints, transports information in the
correct direction, preserves the two initial traces, and remains meaningful
as the null hypersurfaces focus.  This construction comes from the geometric
structure of the equations rather than from a generic fixed-point template.

The numerical implementation combines Legendre--Gauss--Lobatto spectral
elements in the two null coordinates with pole-free spherical
differentiation, angular Galerkin projection, characteristic constraint
solves, and independent overgrid audits.  Extensive AI-assisted software
development made it practical to implement, refactor, and test this large
coupled system.  The mathematical conventions, iteration design, experiment
definitions, acceptance criteria, and scientific interpretation remain
author-controlled; confidence in the computations is based on exact
benchmarks, convergence studies, immutable characteristic data, and
independent residual evaluation rather than on code generation itself.

Eight numerical experiments test complementary aspects of the method.  The
exact-solution tests cover Schwarzschild, Kerr, and Fisher--JNW spacetimes.
They measure the metric error in regular regions, across a Schwarzschild
horizon, and toward a curvature singularity.  The vacuum experiments evolve a
strong nonspherical outgoing perturbation and crossed characteristic data
whose initial curvature is singular at the corner.  The final experiments
evolve nonspherical Einstein--scalar data and a stronger scalar pulse toward
a trapped region.  In the last experiment, an inner atlas of characteristic
rectangles covers a curved domain adapted to null focusing.  Independent
four-metric evaluation finds the section
\begin{equation}\label{eq:introduction-trapped-section}
  (u,v)=(-0.4696271025,0.04)
\end{equation}
on which the two future null expansions are strictly negative everywhere,
with respective spherical suprema $-0.0802152$ and $-4.28011$.  Solving the
angular marginally outer trapped surface equation on successive incoming
null cones then reconstructs an apparent-horizon tube of 18 sections over
$0.00441\leq v\leq0.05$.  Coordinate and angular controls agree on its
$v=0.04$ section to better than $1.6\times10^{-7}$ in the graph location.

This trapped-section computation is motivated by An's proof that anisotropic
outgoing characteristic perturbations of Christodoulou's naked-singularity
data generate an anisotropic apparent horizon that censors the
singularity~\cite{an-25}.  The numerical experiment does not constitute a
proof of that theorem, and the reconstructed tube is tied to the chosen
incoming-null foliation.  It does, however, provide a coordinate- and
angularly controlled numerical realization of the associated trapped-region
mechanism and its anisotropic apparent horizon.

Section~2 fixes the double-null equations and records the analytic framework.
Section~3 constructs the Picard map and its numerical discretization.
Sections~4 and~5 present the vacuum and Einstein--scalar experiments,
respectively, and distinguish exact metric errors from independently
evaluated curvature residuals.
